% Job posting: https://ca.linkedin.com/jobs/view/hearing-performance-developer-at-sonova-group-4463662475
\documentclass[11pt]{article}
\usepackage[left=1in,top=1in,right=1in,bottom=1in]{geometry}
\usepackage[hidelinks]{hyperref}
\usepackage{setspace}
\usepackage[T1]{fontenc}
\usepackage{newtxtext,newtxmath}
\ifdefined\pdfgentounicode
  \input{glyphtounicode}
  \pdfgentounicode=1
\fi
\setstretch{1.1}
\pagenumbering{gobble}
\begin{document}
\pagestyle{empty}

\begin{flushright}
Nishal Stanislaus Silva, PhD \\
Toronto, ON \\
nishal.silva@hotmail.com \\
+1 613 200 2030
\end{flushright}

\noindent Dear Hiring Manager,

\vspace{0.5em}
\noindent As a postdoctoral researcher at the University of Trento, I designed a real-time audio pattern detector that runs in 14 ms end-to-end on a Raspberry Pi 4, holding F1 0.76 while using just 31.4\% CPU -- a 74x speed-up over the 1038 ms baseline it replaced. That result is not a research curiosity; it is evidence that I can take an audio-processing algorithm and make it work inside a razor-thin latency and power budget on constrained silicon, which is precisely the discipline hearing instrument DSP demands: signal-processing decisions that must resolve inside an imperceptible latency window, on a device with almost no power headroom to spare. I am writing to apply for the Hearing Performance Developer role at Sonova.

\vspace{0.5em}
\noindent My toolkit is built for exactly this problem. I work daily with the DSP building blocks that underpin hearing-aid signal chains -- STFT, MFCC/LFCC/PLP/GFCC feature families, the S-transform, CQT and chroma representations, onset and pitch tracking -- and I implement them in C++17 real-time engines (Elk Audio OS, JUCE, VST, OSC) rather than only in offline research code. During a visiting research position at McGill University's IDMIL/CIRMMT lab, I built a self-powered smart electric guitar, fusing an IMU with an audio codec on embedded hardware and profiling compute and power draw to confirm real-time feasibility under a genuinely constrained power budget -- the same battery-constrained mindset that hearing instrument silicon requires. Earlier, as a Research Engineer at MAS Holdings, I spent 5.5 years shipping computer-vision and embedded inference systems into production manufacturing lines, which is where I learned the engineering discipline -- CI/CD, code review, rigorous benchmarking -- that turns a working prototype into a shippable, maintainable product.

\vspace{0.5em}
\noindent I do not have direct industry experience in hearing aids or medical devices, and I want to be upfront about that. What I bring instead is the core transferable skill this role depends on: real-time audio DSP engineering under hard latency and power constraints, validated on embedded hardware and published in peer-reviewed venues. I am confident that gap closes quickly given how closely my technical background maps onto the problem.

\vspace{0.5em}
\noindent I am a Canadian Permanent Resident, eligible to work in Canada without sponsorship, based in Toronto and open to relocating to the Kitchener-Waterloo area for this role. I am available immediately and would welcome the opportunity to discuss how my real-time embedded audio DSP background could contribute to Sonova's hearing performance team.

\vspace{1em}
\noindent Sincerely, \\
Nishal Stanislaus Silva, PhD

\end{document}
